\documentclass{article}
\usepackage{amsmath, amssymb, amsthm}
\usepackage{hyperref}
\usepackage{geometry}
\geometry{margin=1in}

\title{Erd\H{o}s Problem \#106}
\date{Last updated: 2025-08-31}
\author{Source: erdosproblems.com}

\begin{document}
\maketitle

\noindent\textbf{Status:} FALSIFIABLE \quad \textbf{No prize}

\noindent\textbf{Tags:} geometry

\medskip

\noindent\textbf{Problem Statement:}

Draw $n$ squares inside the unit square with no common interior point. Let $f(n)$ be the maximum possible sum of the side-lengths of the squares. Is $f(k^2+1)=k$?




In \cite{Er94b} Erd\H{o}s dates this conjecture to 'more than 60 years ago'. Erd\H{o}s proved that $f(2)=1$ in an early mathematical paper for high school students in Hungary. Newman proved (in personal communication to Erd\H{o}s) that $f(5)=2$. It is trivial from the Cauchy-Schwarz inequality that $f(k^2)=k$. Erd\H{o}s also asks for which $n$ is it true that $f(n+1)=f(n)$. It is easy to see that $f(k^2+1)\geq k$, by first dividing the unit square into $k^2$ smaller squares of side-length $1/k$, and then replacing one square by two smaller squares of side-length $1/2k$. Hal\'{a}sz \cite{Ha84} gives a construction that shows $f(k^2+2)\geq k+\frac{1}{k+1}$, and in general, for any $c\geq 1$,\[f(k^2+2c+1)\geq k+\frac{c}{k}\]and\[f(k^2+2c)\geq k+\frac{c}{k+1}.\]Hal\'{a}sz also considers the variants where we replace a square by a parallelogram or triangle. Erd\H{o}s and Soifer \cite{ErSo95} and Campbell and Staton \cite{CaSt05} have conjectured that, in general, for any integer $-k&#60;c&#60;k$, $f(k^2+2c+1)=k+\frac{c}{k}$, and proved the corresponding lower bound. Praton \cite{Pr08} has proved that this general conjecture is equivalent to $f(k^2+1)=k$. Baek, Koizumi, and Ueoro \cite{BKU24} have proved $g(k^2+1)=k$, where $g(\cdot)$ is defined identically to $f(\cdot)$ with the additional assumption that all squares have sides parallel to the sides of the unit square. More generally, they prove that $g(k^2+2c+1)=k+c/k$ for any $-k&#60;c&#60;k$, which determines all values of $g(\cdot)$. Raj Singh \cite{Ra26} has noted that $f(k^2+1)=k$ being true for all $k$ is equivalent to it being true for infinitely many $k$, which in turn is equivalent to the convergence of\[\sum_{k\geq 1}(f(k^2+1)-k).\]Both equivalences follow from the fact that\[k(f(k^2+1)-k)\]is a non-decreasing function of $k$.




References

[BKU24] Baek, J. and Koizumi, J. and Ueoro, T., A note on the Erd\H{o}s conjecture about square packing . arXiv:2411.07274 (2024).

[CaSt05] Campbell, Connie and Staton, William, A square-packing problem of Erd\H{o}s . Amer. Math. Monthly (2005), 165--167.

[Er94b] Erd\H{o}s, Paul, Some problems in number theory, combinatorics and combinatorial geometry . Math. Pannon. (1994), 261-269.

[ErSo95] Erd\H{o}s, Paul and Soifer, Alexander, Squares in a square . Geombinatorics (1995), 110--114.

[Ha84] Hal\'{a}sz, Sylvia, Packing a convex domain with similar convex domains . J. Combin. Theory Ser. A (1984), 85--90.

[Pr08] Praton, I., Packing squares in a square . Math. Mag. (2008), 358-361.

[Ra26] A. Raj Singh, On a square packing conjecture of Erd\H{o}s . arXiv:2601.22163 (2026).

\bigskip
\noindent\small{Source: \url{https://www.erdosproblems.com/106}}
\end{document}
